% GENERATED FILE. Edit canonical lessons, then run npm run generate:books.
% canonical-source-hash: fnv1a64:b6ff2cdac0f2ee42
% canonical-lessons: TE-C542-citralekhanam, TE-C542-citrincu, TE-C542-silpam, TE-C542-silpi, TE-C542-kalakhandam

\chapter{Painting and Sculpture}
\label{ch:te-painting-and-sculpture}
\hlchaptermodality{542}

\begin{chapteropening}
\textbf{By the end of this chapter:} \emph{I can talk about painting, sculpture and masterpieces.}

Complete the last lesson of chapter 542: I can talk about painting, sculpture and masterpieces.
\end{chapteropening}

\section[\texorpdfstring{\te{చిత్రలేఖనం} (citralēkhanaṁ)}{చిత్రలేఖనం (citralēkhanaṁ)}]{\te{చిత్రలేఖనం} (citralēkhanaṁ) — the art of painting}

\label{lesson:TE-C542-citralekhanam}



\begin{quote}
Before the new one: say the Telugu for a comedian, then the Telugu for spectators.
\end{quote}

\subsection*{You'll want to know: \te{చిత్రలేఖనం}}
\textbf{\te{చిత్రలేఖనం}} — \emph{citralēkhanaṁ} — \textquotedblleft{}the art of painting\textquotedblright{}.

\begin{quote}
\te{నా} \te{చెల్లెలు} \te{చిత్రలేఖనం} \te{నేర్చుకుంటోంది}.
\end{quote}

\begin{quote}
\emph{nā cellelu citralēkhanaṁ nērcukuṇṭōndi.}
\end{quote}

\begin{quote}
My younger sister is learning to paint.
\end{quote}

\te{చిత్రం}, picture, + \te{లేఖనం}, drawing.

\begin{cousinweb}
\textbf{\te{చిత్రలేఖనం}} (\emph{citralēkhanaṁ}) is a learned compound built from Sanskrit \emph{citra}, \textquotedblleft{}bright, variegated\textquotedblright{} (CDIAL 4803), and \emph{lēkhana}, \textquotedblleft{}scraping\textquotedblright{}, the noun of the root of writing (CDIAL 11103).
\end{cousinweb}

\begin{grammarlens}[title={What you've built}]
One more word for describing what you lived and felt.
\end{grammarlens}

\subsection*{Your turn}
\begin{itemize}
\raggedright
  \item \emph{Say it:} \emph{citralēkhanaṁ}
  \item \emph{Say it:} \emph{citralēkhanaṁ}, once more
  \item \emph{Say it:} \emph{nā cellelu citralēkhanaṁ nērcukuṇṭōndi.} — My younger sister is learning to paint.
  \item \emph{Recall:} say the Telugu for a comedian, then the Telugu for spectators, then say \emph{citralēkhanaṁ} again
\end{itemize}

\begin{tcolorbox}[breakable,colback=teal!4,colframe=teal!35!black,title={Before you move on}]
What does \te{చిత్రలేఖనం} mean? (\textquotedblleft{}The art of painting\textquotedblright{}.) Say it once more.
\end{tcolorbox}
\section[\texorpdfstring{\te{చిత్రించు} (citriñcu)}{చిత్రించు (citriñcu)}]{\te{చిత్రించు} (citriñcu) — to portray; to paint}

\label{lesson:TE-C542-citrincu}



\begin{quote}
Before the new one: say the Telugu for spectators, then the Telugu for the art of painting.
\end{quote}

\subsection*{You'll want to know: \te{చిత్రించు}}
\textbf{\te{చిత్రించు}} — \emph{citriñcu} — \textquotedblleft{}to portray; to paint\textquotedblright{}.

\begin{quote}
\te{ఆ} \te{కళాకారుడు} \te{పల్లె} \te{జీవితాన్ని} \te{చిత్రించాడు}.
\end{quote}

\begin{quote}
\emph{ā kaḷākāruḍu palle jīvitānni citriñcāḍu.}
\end{quote}

\begin{quote}
That artist portrayed village life.
\end{quote}

Past (I …): \te{చిత్రించాను}, \emph{citriñcānu}.

In paint or in words.

\begin{cousinweb}
\textbf{\te{చిత్రించు}} (\emph{citriñcu}) is built on Sanskrit \emph{citra}, \textquotedblleft{}bright, variegated\textquotedblright{}, the source of \te{చిత్రం}, \textquotedblleft{}a picture\textquotedblright{} (CDIAL 4803), with the Telugu verb ending -\te{ఇంచు}.
\end{cousinweb}

\begin{grammarlens}[title={What you've built}]
One more word for describing what you lived and felt.
\end{grammarlens}

\subsection*{Your turn}
\begin{itemize}
\raggedright
  \item \emph{Say it:} \emph{citriñcu}
  \item \emph{Say it:} \emph{citriñcu}, once more
  \item \emph{Say it:} \emph{ā kaḷākāruḍu palle jīvitānni citriñcāḍu.} — That artist portrayed village life.
  \item \emph{Recall:} say the Telugu for spectators, then the Telugu for the art of painting, then say \emph{citriñcu} again
\end{itemize}

\begin{tcolorbox}[breakable,colback=teal!4,colframe=teal!35!black,title={Before you move on}]
What does \te{చిత్రించు} mean? (\textquotedblleft{}To portray; to paint\textquotedblright{}.) Say it once more.
\end{tcolorbox}
\section[\texorpdfstring{\te{శిల్పం} (śilpaṁ)}{శిల్పం (śilpaṁ)}]{\te{శిల్పం} (śilpaṁ) — a sculpture; a carving}

\label{lesson:TE-C542-silpam}



\begin{quote}
Before the new one: say the Telugu for the art of painting, then the Telugu for to portray.
\end{quote}

\subsection*{You'll want to know: \te{శిల్పం}}
\textbf{\te{శిల్పం}} — \emph{śilpaṁ} — \textquotedblleft{}a sculpture; a carving\textquotedblright{}.

\begin{quote}
\te{ఆ} \te{గుడిలో} \te{అందమైన} \te{శిల్పాలు} \te{ఉన్నాయి}.
\end{quote}

\begin{quote}
\emph{ā guḍilō andamaina śilpālu unnāyi.}
\end{quote}

\begin{quote}
There are beautiful sculptures in that temple.
\end{quote}

Plural: \te{శిల్పాలు}, \emph{śilpālu}.

\te{శిల్పకళ}, the art of sculpture.

\begin{cousinweb}
\textbf{\te{శిల్పం}} (\emph{śilpaṁ}) is borrowed from Sanskrit \emph{śilpa}, \textquotedblleft{}artistic work; manual craft\textquotedblright{} (CDIAL 12470).
\end{cousinweb}

\begin{grammarlens}[title={What you've built}]
One more word for describing what you lived and felt.
\end{grammarlens}

\subsection*{Your turn}
\begin{itemize}
\raggedright
  \item \emph{Say it:} \emph{śilpaṁ}
  \item \emph{Say it:} \emph{śilpaṁ}, once more
  \item \emph{Say it:} \emph{ā guḍilō andamaina śilpālu unnāyi.} — There are beautiful sculptures in that temple.
  \item \emph{Recall:} say the Telugu for the art of painting, then the Telugu for to portray, then say \emph{śilpaṁ} again
\end{itemize}

\begin{tcolorbox}[breakable,colback=teal!4,colframe=teal!35!black,title={Before you move on}]
What does \te{శిల్పం} mean? (\textquotedblleft{}A sculpture; a carving\textquotedblright{}.) Say it once more.
\end{tcolorbox}
\section[\texorpdfstring{\te{శిల్పి} (śilpi)}{శిల్పి (śilpi)}]{\te{శిల్పి} (śilpi) — a sculptor}

\label{lesson:TE-C542-silpi}



\begin{quote}
Before the new one: say the Telugu for to portray, then the Telugu for a sculpture.
\end{quote}

\subsection*{You'll want to know: \te{శిల్పి}}
\textbf{\te{శిల్పి}} — \emph{śilpi} — \textquotedblleft{}a sculptor\textquotedblright{}.

\begin{quote}
\te{ఆ} \te{శిల్పి} \te{రాతితో} \te{బొమ్మలు} \te{చెక్కుతాడు}.
\end{quote}

\begin{quote}
\emph{ā śilpi rātitō bommalu cekkutāḍu.}
\end{quote}

\begin{quote}
That sculptor carves figures from stone.
\end{quote}

Plural \te{శిల్పులు}. \te{చెక్కు}, to carve.

\begin{cousinweb}
\textbf{\te{శిల్పి}} (\emph{śilpi}) is borrowed from Sanskrit \emph{śilpin}, \textquotedblleft{}skilled in art; an artificer\textquotedblright{} (CDIAL 12471).
\end{cousinweb}

\begin{grammarlens}[title={What you've built}]
One more word for describing what you lived and felt.
\end{grammarlens}

\subsection*{Your turn}
\begin{itemize}
\raggedright
  \item \emph{Say it:} \emph{śilpi}
  \item \emph{Say it:} \emph{śilpi}, once more
  \item \emph{Say it:} \emph{ā śilpi rātitō bommalu cekkutāḍu.} — That sculptor carves figures from stone.
  \item \emph{Recall:} say the Telugu for to portray, then the Telugu for a sculpture, then say \emph{śilpi} again
\end{itemize}

\begin{tcolorbox}[breakable,colback=teal!4,colframe=teal!35!black,title={Before you move on}]
What does \te{శిల్పి} mean? (\textquotedblleft{}A sculptor\textquotedblright{}.) Say it once more.
\end{tcolorbox}
\section[\texorpdfstring{\te{కళాఖండం} (kaḷākhaṇḍaṁ)}{కళాఖండం (kaḷākhaṇḍaṁ)}]{\te{కళాఖండం} (kaḷākhaṇḍaṁ) — a masterpiece}

\label{lesson:TE-C542-kalakhandam}



\begin{quote}
Before the new one: say the Telugu for a sculpture, then the Telugu for a sculptor.
\end{quote}

\subsection*{You'll want to know: \te{కళాఖండం}}
\textbf{\te{కళాఖండం}} — \emph{kaḷākhaṇḍaṁ} — \textquotedblleft{}a masterpiece\textquotedblright{}.

\begin{quote}
\te{ఆ} \te{పాత} \te{ఆలయం} \te{ఒక} \te{కళాఖండం}.
\end{quote}

\begin{quote}
\emph{ā pāta ālayaṁ oka kaḷākhaṇḍaṁ.}
\end{quote}

\begin{quote}
That old temple is a masterpiece.
\end{quote}

\te{కళ}, art, + \te{ఖండం}, piece.

\begin{cousinweb}
\textbf{\te{కళాఖండం}} (\emph{kaḷākhaṇḍaṁ}) is a learned compound built from Sanskrit \emph{kalā}, \textquotedblleft{}any practical art\textquotedblright{} (CDIAL 2927), and \emph{khaṇḍa}, \textquotedblleft{}a piece, a fragment\textquotedblright{} (CDIAL 3792).
\end{cousinweb}

\begin{grammarlens}[title={What you've built}]
One more word for describing what you lived and felt.
\end{grammarlens}

\subsection*{Your turn}
\begin{itemize}
\raggedright
  \item \emph{Say it:} \emph{kaḷākhaṇḍaṁ}
  \item \emph{Say it:} \emph{kaḷākhaṇḍaṁ}, once more
  \item \emph{Say it:} \emph{ā pāta ālayaṁ oka kaḷākhaṇḍaṁ.} — That old temple is a masterpiece.
  \item \emph{Recall:} say the Telugu for a sculpture, then the Telugu for a sculptor, then say \emph{kaḷākhaṇḍaṁ} again
\end{itemize}

\begin{tcolorbox}[breakable,colback=teal!4,colframe=teal!35!black,title={Before you move on}]
What does \te{కళాఖండం} mean? (\textquotedblleft{}A masterpiece\textquotedblright{}.) Say it once more.
\end{tcolorbox}
